%%% =====================================================================
%%% macros.tex --- CQ-automata project (stemmed LSTA + classical-quantum
%%%                program verification)
%%%
%%% Single-file macros for this project.  Section layout mirrors the
%%% modular layout in JAC's latex-macros repository
%%% (https://.../latex-macros/macros/{core,automata,verification}.tex)
%%% so that pieces can be migrated out into the shared library later.
%%%
%%% Loaded once via   %%% =====================================================================
%%% macros.tex --- CQ-automata project (stemmed LSTA + classical-quantum
%%%                program verification)
%%%
%%% Single-file macros for this project.  Section layout mirrors the
%%% modular layout in JAC's latex-macros repository
%%% (https://.../latex-macros/macros/{core,automata,verification}.tex)
%%% so that pieces can be migrated out into the shared library later.
%%%
%%% Loaded once via   %%% =====================================================================
%%% macros.tex --- CQ-automata project (stemmed LSTA + classical-quantum
%%%                program verification)
%%%
%%% Single-file macros for this project.  Section layout mirrors the
%%% modular layout in JAC's latex-macros repository
%%% (https://.../latex-macros/macros/{core,automata,verification}.tex)
%%% so that pieces can be migrated out into the shared library later.
%%%
%%% Loaded once via   %%% =====================================================================
%%% macros.tex --- CQ-automata project (stemmed LSTA + classical-quantum
%%%                program verification)
%%%
%%% Single-file macros for this project.  Section layout mirrors the
%%% modular layout in JAC's latex-macros repository
%%% (https://.../latex-macros/macros/{core,automata,verification}.tex)
%%% so that pieces can be migrated out into the shared library later.
%%%
%%% Loaded once via   \input{macros}   in main.tex.
%%%
%%% Conventions
%%%   - Sets in \mathbb           (\nat, \natz, \bools, \complex)
%%%   - Bit alphabet              \bit = \{0,1\}
%%%   - Tuples                    \tuple{a, b}
%%%   - Tool names                \autoq, \autq, \tool, ...
%%%   - LSTA / sLSTA names        \aut = \mathcal{A}; root states \rootstates
%%%   - Stem alphabet alias       \Store = Sigma_1 (rank-1 stem labels)
%%%   - Author notes              \ol, \yfc, \jal, \wl, \fil
%%%
%%% Sections
%%%   1.  Author notes and dev macros
%%%   2.  Sets, numbers, words
%%%   3.  Tree primitives (perfect binary, stemmed perfect binary)
%%%   4.  LSTA / sLSTA: states, transitions, runs, languages
%%%   5.  Transition-arrow rendering helpers
%%%   6.  Quantum gates and operations
%%%   7.  Specifications and Hoare-style triples
%%%   8.  Tool names and method acronyms
%%%   9.  Small tree-drawing helpers
%%%  10.  Miscellaneous
%%% =====================================================================


%%% =====================================================================
%%% 1. Author notes and dev macros
%%% =====================================================================

%%% Inline coloured author notes.  Toggle off by redefining to \gobble.
\newcommand{\ol}[1]{\textcolor{blue}{\ifmmode \text{[OL: #1]}\else [OL: #1] \fi}}
\newcommand{\yfc}[1]{\textcolor{green!50!blue}{\ifmmode \text{[YF: #1]}\else [YF: #1] \fi}}
\newcommand{\jal}[1]{\textcolor{pink}{\ifmmode \text{[JA: #1]}\else [JA: #1] \fi}}
\newcommand{\wl}[1]{\textcolor{orange}{\ifmmode \text{[WL: #1]}\else [WL: #1] \fi}}
\newcommand{\fil}[1]{\textcolor{red}{\ifmmode \text{[FI: #1]}\else [FI: #1] \fi}}

%%% Hide content (handy for parking paragraphs).  \showit kept for symmetry.
\newcommand{\hide}[1]{}
\newcommand{\showit}[1]{#1}

%%% Marker for places where a definition was added during a rewrite.
%%% Currently no-op; switch to the commented version to highlight.
% \newcommand{\newdef}[1]{\textcolor{purple}{[newdef: #1]}}
\newcommand{\newdef}[1]{}


%%% =====================================================================
%%% 2. Sets, numbers, words
%%% =====================================================================

\newcommand{\nat}{\mathbb{N}}                 % positive naturals (no 0)
\newcommand{\natz}{\nat_{0}}                  % naturals with 0
\newcommand{\naturals}{\nat}                  % alias
\newcommand{\integers}{\mathbb{Z}}
\newcommand{\bools}{\mathbb{B}}
\newcommand{\complex}{\mathbb{C}}
\newcommand{\cnum}{\complex}                  % older alias kept for back-compat

%%% Alphabets and bit strings
\newcommand{\bit}{\{0,1\}}                    % binary alphabet
\newcommand{\unary}{\{\mathtt{1}\}}           % unary alphabet (legacy)
\newcommand{\abc}{\Gamma}                     % generic alphabet
\newcommand{\colors}{\Omega}                  % colour / choice universe (legacy)

%%% The empty word.  Sometimes typeset epsilon, sometimes varepsilon; we
%%% standardise on \varepsilon for readability.
\newcommand{\emptyword}{\varepsilon}

%%% Tuples and scalable angle brackets.
%%% \tuple is a plain macro (matches existing usage q = \tuple{Q, ...}).
\newcommand{\tuple}[1]{\langle #1 \rangle}

%%% Common short fractions used in amplitudes.
\newcommand{\smallfrac}[2]{{\textstyle \frac{#1}{#2}}}
\newcommand{\invtwo}{\smallfrac{1}{2}}
\newcommand{\invsqrttwo}{\smallfrac{1}{\sqrt{2}}}


%%% =====================================================================
%%% 3. Tree primitives
%%%
%%% Covers both the perfect binary trees of [LSTA2025] and the
%%% stemmed perfect binary trees introduced in this paper.
%%% =====================================================================

%%% ---- generic tree language and domain --------------------------------
\newcommand{\trees}{\mathbb{T}}               % set of all (ranked) trees
\newcommand{\treesof}[1]{\trees_{#1}}
\newcommand{\abctrees}{\treesof{\abc}}

\newcommand{\dom}{\mathrm{dom}}               % \dom(T) -- domain of a tree
\newcommand{\domof}[1]{\dom(#1)}
\newcommand{\height}[1]{h(#1)}                % alternate height notation
\newcommand{\htof}[1]{\mathrm{ht}(#1)}        % preferred: \htof{v}

%%% Subtree, tree-cons, dummy zero tree, concatenation.
\newcommand{\subtreeof}[2]{#1_{|#2}}
\newcommand{\treecons}[3]{\mathsf{cons}(#1, #2, #3)}
\newcommand{\tzeroof}[1]{t^{#1}_{\mathit{zero}}}
\newcommand{\concat}{\mathbin{\bullet}}

%%% ---- ranked alphabets ------------------------------------------------
%%% Arity function on a ranked alphabet.  \arity{\Sigma} -> '#'.
%%% Defined as a macro so we can refer to it both in display ($\arity\colon
%%% \Sigma \to \{0,1,2\}$) and in prose ("\arity-1 symbols").
\newcommand{\arity}{\#}
\newcommand{\arityof}[1]{\arity(#1)}

%%% Stem-alphabet alias used by the sLSTA: Store = Sigma_1, the rank-1
%%% symbols carrying classical-store information.
\newcommand{\Store}{\mathtt{Store}}

%%% ---- shape of a stemmed perfect binary tree --------------------------
%%% (m, n) = (stem length, binary height).  Use \shape{m}{n}.
\newcommand{\shape}[2]{(#1, #2)}

%%% Root of a (possibly stemmed) tree T.  v_star with a star subscript.
\newcommand{\rootnode}{v_{\star}}


%%% =====================================================================
%%% 4. LSTA / sLSTA: states, transitions, runs, languages
%%% =====================================================================

%%% ---- automata names (calligraphic) ----------------------------------
\newcommand{\aut}{\mathcal{A}}
\newcommand{\autb}{\mathcal{B}}
\newcommand{\autm}{\mathcal{M}}
\newcommand{\autinv}{\mathcal{I}}
\newcommand{\cut}{\mathcal{C}}                % combined / context automaton
\newcommand{\nut}{\mathcal{N}}
\newcommand{\autex}{\aut_{\mathit{ex}}}       % a running-example automaton
\newcommand{\butex}{\autb_{\mathit{ex}}}
\newcommand{\autpre}{\aut_{\mathrm{pre}}}
\newcommand{\autpost}{\aut_{\mathrm{post}}}
\newcommand{\autbases}{\aut_{\mathit{bases}}}
\newcommand{\autRes}{\aut_{\mathsf{Result}}}
\newcommand{\auto}{\mathscr{A}}               % script-A (rarely used)

%%% ---- root states, transitions ----------------------------------------
\newcommand{\rootstates}{\mathcal{R}}

%%% A transition rendered as   q  --[C]-->  body.
%%% Body is either a binary symbol f(q0,q1), a unary symbol b(p), or a
%%% leaf symbol f.
\newcommand{\transtree}[3]{{#1}\xrightarrow{#2}~{#3}}

%%% Selectors on a transition delta.
\newcommand{\topof}[1]{\mathtt{top}(#1)}      % source state
\newcommand{\symof}[1]{\mathtt{sym}(#1)}      % label symbol
\newcommand{\botof}[1]{\mathtt{bot}(#1)}      % target state(s)

%%% ---- runs and level synchronization ---------------------------------
\newcommand{\run}{\rho}                       % canonical run variable
\newcommand{\level}{\mathrm{level}}           % \level(\run, d)
\newcommand{\choice}{\mathrm{choice}}         % \choice(\run)

%%% Language and choice-indexed language.
\newcommand{\lang}{\mathcal{L}}
\newcommand{\langof}[1]{\lang(#1)}

%%% Semantic brackets [[ . ]].  Provided by stmaryrd; we fall back to
%%% ASCII look-alikes if the package is not available.
\IfFileExists{stmaryrd.sty}{\usepackage{stmaryrd}}{%
  \providecommand{\llbracket}{[\![}%
  \providecommand{\rrbracket}{]\!]}%
}
\newcommand{\semof}[1]{\llbracket #1 \rrbracket}
\newcommand{\semofof}[2]{\semof{#1}(#2)}

%%% Up-to-scale entailment relation used for measurement-up-to-norm.
\newcommand{\entailedbyuptosc}{\mathrel{\models^{\mathit{uts}}}}


%%% =====================================================================
%%% 5. Transition-arrow rendering helpers
%%%
%%% Used in figures and inline transition diagrams.  Kept as-is from the
%%% legacy macros; they construct compact labelled arrows of the form
%%%       --[label]-->
%%% =====================================================================

\makeatletter
\DeclareRobustCommand{\shortto}{%
  \mathrel{\mathpalette\short@to\relax}%
}
\DeclareRobustCommand{\shortminus}{%
  \mathrel{\mathpalette\short@minus\relax}%
}
\newcommand{\short@to}[2]{%
  \mkern2mu
  \clipbox{{.5\width} 0 0 0}{$\m@th#1\vphantom{+}{\rightarrow}$}%
}
\newcommand{\short@minus}[2]{%
  \mkern2mu
  \clipbox{{.5\width} 0 0 0}{$\m@th#1\vphantom{+}{-}$}%
}
\makeatother

\newcommand{\labeledto}[1]{%
  \raisebox{-0.2pt}{\scalebox{1.2}{\ensuremath{%
    {\shortminus}\hspace{-1.7pt}%
    \raisebox{0.16ex}{$\scriptstyle[#1\hspace{-0.28pt}]$}%
    \hspace{-2.25pt}{\shortto}}}}}
\newcommand{\newlabeledto}[1]{%
  \mbox{$\longrightarrow\rlap{\hspace*{-4.95mm}%
    \raisebox{0.00mm}{\scalebox{0.8}{#1}}}$\!}}

\newcommand{\move}[1]{\labeledto{#1}}

%%% LTS-style transitions (used in transition systems, not LSTA proper).
\newcommand{\ltstr}{\mathit{Tr}}
\newcommand{\ltr}[1]{\xrightarrow{\smash{\raisebox{-0.7mm}{$\scriptstyle #1$}}}}

%%% Named-form transition macros from the legacy file (kept for the few
%%% places that build figures with them).
\newcommand{\trntransof}[4]{#1 \to #2(#3, #4)}
\newcommand{\transof}[5]{#1 \newlabeledto{#3} #2(#4, #5)}
\newcommand{\trnimageof}[2]{#1(#2)}
\newcommand{\trnimagestateof}[3]{\trnimageof{#1}{#2, #3}}


%%% =====================================================================
%%% 6. Quantum gates and operations
%%% =====================================================================

\newcommand{\ketof}[1]{\ket{#1}}

%%% Gate "wrappers" used in math mode.
\newcommand{\gateof}[1]{\mathrm{#1}}
\newcommand{\matrixof}[1]{\mathit{mat}(#1)}
\newcommand{\rowof}[1]{\mathit{row}(#1)}
\newcommand{\vectorof}[1]{\mathit{vec}(#1)}

%%% Common single- and two-qubit gates.
\newcommand{\had}{\mathsf{H}}
\newcommand{\hadtensor}{\had^{\otimes}}
\newcommand{\sgate}{\mathsf{S}}
\newcommand{\xgate}{\mathsf{X}}
\newcommand{\idgate}{\mathrm{Id}}
\newcommand{\cnot}{\mathsf{CX}}
\newcommand{\ccnot}{\mathsf{CCX}}
\newcommand{\rotx}{\mathsf{R}_{x}(\tfrac{\pi}{2})}

%%% Measurement.   \measurement{i}{=b}  -- measure qubit i, observe b.
\newcommand{\measurement}[2]{M_{#1}{#2}}

%%% Tree-transformation helpers (single-gate construction labels L/R).
\newcommand{\leftT}{\mathtt{L}}
\newcommand{\rightT}{\mathtt{R}}

%%% Classical store and qubit "directions" used in sLSTA constructions.
%%% These name the two principal directions of a classical branch.
\newcommand{\leftof}[1]{\mathtt{left}(#1)}
\newcommand{\rightof}[1]{\mathtt{right}(#1)}


%%% =====================================================================
%%% 7. Specifications and Hoare-style triples
%%% =====================================================================

%%% Standard Hoare triple {P} c {Q}.
\newcommand{\hoaretrip}[3]{\{#1\}~{#2}~\{#3\}}
%%% "Functional" triple with brackets (|P|) c (|Q|), used for specs
%%% intended to encode a Choi--Jamiolkowski-style state pair.
\newcommand{\assert}[1]{(\!|{#1}|\!)}
\newcommand{\funtrip}[3]{\assert{#1}\; {#2}\; \assert{#3}}

%%% Spec pair  c | psi >  used to bind a star pre/post to a state.
\newcommand{\specpair}[2]{#2\ket{#1}}
\newcommand{\specdef}[1]{#1\specpair{*}{}}

%%% Pre-/post-condition standard names (kept for cross-paper consistency).
\newcommand{\PreCond}{P}
\newcommand{\PostCond}{Q}
\newcommand{\pre}{\mathrm{Pre}}
\newcommand{\post}{\mathrm{Post}}

%%% Vectors of all-zeros / all-ones; the imaginary-zero subspace name.
\newcommand{\veczero}{\vec{0}}
\newcommand{\vecone}{\vec{1}}
\newcommand{\imzero}{\psi_{\Im}}


%%% =====================================================================
%%% 8. Tool names and method acronyms
%%%
%%% \tool refers to *our* tool; aliases below are for related-work and
%%% comparison tables.
%%% =====================================================================

\newcommand{\tool}{\textsc{CQ-AutoQ}\xspace}    % working name -- adjust later
\newcommand{\autoq}{\textsc{AutoQ}\xspace}
\newcommand{\sliqsim}{\textsc{SliQSim}\xspace}
\newcommand{\svsim}{\textsc{SV-Sim}\xspace}
\newcommand{\sliqec}{\textsc{SliQEC}\xspace}
\newcommand{\feynman}{\textsc{Feynman}\xspace}
\newcommand{\feynopt}{\textsc{Feynopt}\xspace}
\newcommand{\qcec}{\textsc{Qcec}\xspace}
\newcommand{\quasimodo}{\textsc{Quasimodo}\xspace}
\newcommand{\qbricks}{\textsc{Qbricks}\xspace}
\newcommand{\qiskit}{\textsc{Qiskit}\xspace}
\newcommand{\vata}{\textsc{Vata}\xspace}
\newcommand{\quartz}{\textsc{Quartz}\xspace}
\newcommand{\symqv}{\textsc{symQV}\xspace}
\newcommand{\caal}{\textsc{CaAL}\xspace}
\newcommand{\isabelle}{\textsc{Isabelle}\xspace}
\newcommand{\coq}{\textsc{Coq}\xspace}
\newcommand{\qasm}{\textsc{OpenQASM}\xspace}

%%% Method acronyms.
\newcommand{\lsta}{LSTA\xspace}
\newcommand{\lstas}{LSTAs\xspace}
\newcommand{\slsta}{sLSTA\xspace}
\newcommand{\slstas}{sLSTAs\xspace}

%%% Complexity classes (lower-case to avoid clash with cleveref labels).
\newcommand{\pspace}{\textsc{PSpace}\xspace}
\newcommand{\expspace}{\textsc{ExpSpace}\xspace}
\newcommand{\bigO}{\mathcal{O}}


%%% =====================================================================
%%% 9. Small tree-drawing helpers
%%%
%%% Inline TikZ snippets for compact tree literals in prose.  Pulled
%%% from the legacy file unchanged.
%%% =====================================================================

\newcommand{\smalltreeof}[3]{%
\scalebox{0.75}{%
\raisebox{3mm}{%
\tikz[level distance=2mm,sibling distance=8mm,baseline,inner sep=0.2mm]{%
\node {$#1$} child {node {$#2$}}%
child {node {$#3$}};}}}%
}

\newcommand{\bigtreeof}[6]{%
\scalebox{0.75}{%
\raisebox{3mm}{%
\tikz[level distance=2mm,
  level 1/.style={sibling distance=12mm},
  level 2/.style={sibling distance=8mm},
  baseline,inner sep=0.2mm]{%
\node {$#1$} child {node {$#2$}
                    child {node {$#3$}}
                    child {node {$#4$}}
                   }%
             child {node {$#2$}
                    child {node {$#5$}}
                    child {node {$#6$}}
                   };
}}}%
}


%%% =====================================================================
%%% 10. Miscellaneous
%%% =====================================================================

%%% Support set of a linear / weighted combination.
\newcommand{\supp}{\text{st}}
\newcommand{\suppof}[1]{\supp(#1)}

%%% Linear-form notation used in some proofs.
\newcommand{\linforms}{\mathbb{L}}
\newcommand{\linformsof}[1]{\linforms_{#1}}
\newcommand{\linof}[1]{\boldsymbol{#1}}
\newcommand{\linindexof}[2]{\linof{#1}[#2]}

%%% Positive part [x]_+.
\newcommand{\mm}[1]{[#1]_{+}}

%%% Misc symbols carried over from prior projects.
\newcommand{\sep}{\texttt{\#}}                % separator symbol
\newcommand{\initmark}{\to}
\newcommand{\partialto}{\rightharpoonup}
\newcommand{\compose}{\mathbin{\circ}}
\newcommand{\wildcard}{\mathtt{?}}

%%% Counting / encoding helpers.
\newcommand{\numof}[1]{\mathit{num}(#1)}
\newcommand{\revof}[1]{\mathit{rev}(#1)}
\newcommand{\encof}[1]{\mathit{enc}(#1)}

%%% A generic arithmetic operator placeholder used in some proofs.
\newcommand{\arithop}{\mathbin{\circledast}}
\newcommand{\minus}{\mathit{minus}}

%%% Choice / colour strings (legacy).
\newcommand{\colorstr}[1]{\mathbf{colorstr}(#1)}
\newcommand{\coloraut}[1]{\mathbf{dom}(#1)}

%%% ---- end of macros.tex ---------------------------------------------
   in main.tex.
%%%
%%% Conventions
%%%   - Sets in \mathbb           (\nat, \natz, \bools, \complex)
%%%   - Bit alphabet              \bit = \{0,1\}
%%%   - Tuples                    \tuple{a, b}
%%%   - Tool names                \autoq, \autq, \tool, ...
%%%   - LSTA / sLSTA names        \aut = \mathcal{A}; root states \rootstates
%%%   - Stem alphabet alias       \Store = Sigma_1 (rank-1 stem labels)
%%%   - Author notes              \ol, \yfc, \jal, \wl, \fil
%%%
%%% Sections
%%%   1.  Author notes and dev macros
%%%   2.  Sets, numbers, words
%%%   3.  Tree primitives (perfect binary, stemmed perfect binary)
%%%   4.  LSTA / sLSTA: states, transitions, runs, languages
%%%   5.  Transition-arrow rendering helpers
%%%   6.  Quantum gates and operations
%%%   7.  Specifications and Hoare-style triples
%%%   8.  Tool names and method acronyms
%%%   9.  Small tree-drawing helpers
%%%  10.  Miscellaneous
%%% =====================================================================


%%% =====================================================================
%%% 1. Author notes and dev macros
%%% =====================================================================

%%% Inline coloured author notes.  Toggle off by redefining to \gobble.
\newcommand{\ol}[1]{\textcolor{blue}{\ifmmode \text{[OL: #1]}\else [OL: #1] \fi}}
\newcommand{\yfc}[1]{\textcolor{green!50!blue}{\ifmmode \text{[YF: #1]}\else [YF: #1] \fi}}
\newcommand{\jal}[1]{\textcolor{pink}{\ifmmode \text{[JA: #1]}\else [JA: #1] \fi}}
\newcommand{\wl}[1]{\textcolor{orange}{\ifmmode \text{[WL: #1]}\else [WL: #1] \fi}}
\newcommand{\fil}[1]{\textcolor{red}{\ifmmode \text{[FI: #1]}\else [FI: #1] \fi}}

%%% Hide content (handy for parking paragraphs).  \showit kept for symmetry.
\newcommand{\hide}[1]{}
\newcommand{\showit}[1]{#1}

%%% Marker for places where a definition was added during a rewrite.
%%% Currently no-op; switch to the commented version to highlight.
% \newcommand{\newdef}[1]{\textcolor{purple}{[newdef: #1]}}
\newcommand{\newdef}[1]{}


%%% =====================================================================
%%% 2. Sets, numbers, words
%%% =====================================================================

\newcommand{\nat}{\mathbb{N}}                 % positive naturals (no 0)
\newcommand{\natz}{\nat_{0}}                  % naturals with 0
\newcommand{\naturals}{\nat}                  % alias
\newcommand{\integers}{\mathbb{Z}}
\newcommand{\bools}{\mathbb{B}}
\newcommand{\complex}{\mathbb{C}}
\newcommand{\cnum}{\complex}                  % older alias kept for back-compat

%%% Alphabets and bit strings
\newcommand{\bit}{\{0,1\}}                    % binary alphabet
\newcommand{\unary}{\{\mathtt{1}\}}           % unary alphabet (legacy)
\newcommand{\abc}{\Gamma}                     % generic alphabet
\newcommand{\colors}{\Omega}                  % colour / choice universe (legacy)

%%% The empty word.  Sometimes typeset epsilon, sometimes varepsilon; we
%%% standardise on \varepsilon for readability.
\newcommand{\emptyword}{\varepsilon}

%%% Tuples and scalable angle brackets.
%%% \tuple is a plain macro (matches existing usage q = \tuple{Q, ...}).
\newcommand{\tuple}[1]{\langle #1 \rangle}

%%% Common short fractions used in amplitudes.
\newcommand{\smallfrac}[2]{{\textstyle \frac{#1}{#2}}}
\newcommand{\invtwo}{\smallfrac{1}{2}}
\newcommand{\invsqrttwo}{\smallfrac{1}{\sqrt{2}}}


%%% =====================================================================
%%% 3. Tree primitives
%%%
%%% Covers both the perfect binary trees of [LSTA2025] and the
%%% stemmed perfect binary trees introduced in this paper.
%%% =====================================================================

%%% ---- generic tree language and domain --------------------------------
\newcommand{\trees}{\mathbb{T}}               % set of all (ranked) trees
\newcommand{\treesof}[1]{\trees_{#1}}
\newcommand{\abctrees}{\treesof{\abc}}

\newcommand{\dom}{\mathrm{dom}}               % \dom(T) -- domain of a tree
\newcommand{\domof}[1]{\dom(#1)}
\newcommand{\height}[1]{h(#1)}                % alternate height notation
\newcommand{\htof}[1]{\mathrm{ht}(#1)}        % preferred: \htof{v}

%%% Subtree, tree-cons, dummy zero tree, concatenation.
\newcommand{\subtreeof}[2]{#1_{|#2}}
\newcommand{\treecons}[3]{\mathsf{cons}(#1, #2, #3)}
\newcommand{\tzeroof}[1]{t^{#1}_{\mathit{zero}}}
\newcommand{\concat}{\mathbin{\bullet}}

%%% ---- ranked alphabets ------------------------------------------------
%%% Arity function on a ranked alphabet.  \arity{\Sigma} -> '#'.
%%% Defined as a macro so we can refer to it both in display ($\arity\colon
%%% \Sigma \to \{0,1,2\}$) and in prose ("\arity-1 symbols").
\newcommand{\arity}{\#}
\newcommand{\arityof}[1]{\arity(#1)}

%%% Stem-alphabet alias used by the sLSTA: Store = Sigma_1, the rank-1
%%% symbols carrying classical-store information.
\newcommand{\Store}{\mathtt{Store}}

%%% ---- shape of a stemmed perfect binary tree --------------------------
%%% (m, n) = (stem length, binary height).  Use \shape{m}{n}.
\newcommand{\shape}[2]{(#1, #2)}

%%% Root of a (possibly stemmed) tree T.  v_star with a star subscript.
\newcommand{\rootnode}{v_{\star}}


%%% =====================================================================
%%% 4. LSTA / sLSTA: states, transitions, runs, languages
%%% =====================================================================

%%% ---- automata names (calligraphic) ----------------------------------
\newcommand{\aut}{\mathcal{A}}
\newcommand{\autb}{\mathcal{B}}
\newcommand{\autm}{\mathcal{M}}
\newcommand{\autinv}{\mathcal{I}}
\newcommand{\cut}{\mathcal{C}}                % combined / context automaton
\newcommand{\nut}{\mathcal{N}}
\newcommand{\autex}{\aut_{\mathit{ex}}}       % a running-example automaton
\newcommand{\butex}{\autb_{\mathit{ex}}}
\newcommand{\autpre}{\aut_{\mathrm{pre}}}
\newcommand{\autpost}{\aut_{\mathrm{post}}}
\newcommand{\autbases}{\aut_{\mathit{bases}}}
\newcommand{\autRes}{\aut_{\mathsf{Result}}}
\newcommand{\auto}{\mathscr{A}}               % script-A (rarely used)

%%% ---- root states, transitions ----------------------------------------
\newcommand{\rootstates}{\mathcal{R}}

%%% A transition rendered as   q  --[C]-->  body.
%%% Body is either a binary symbol f(q0,q1), a unary symbol b(p), or a
%%% leaf symbol f.
\newcommand{\transtree}[3]{{#1}\xrightarrow{#2}~{#3}}

%%% Selectors on a transition delta.
\newcommand{\topof}[1]{\mathtt{top}(#1)}      % source state
\newcommand{\symof}[1]{\mathtt{sym}(#1)}      % label symbol
\newcommand{\botof}[1]{\mathtt{bot}(#1)}      % target state(s)

%%% ---- runs and level synchronization ---------------------------------
\newcommand{\run}{\rho}                       % canonical run variable
\newcommand{\level}{\mathrm{level}}           % \level(\run, d)
\newcommand{\choice}{\mathrm{choice}}         % \choice(\run)

%%% Language and choice-indexed language.
\newcommand{\lang}{\mathcal{L}}
\newcommand{\langof}[1]{\lang(#1)}

%%% Semantic brackets [[ . ]].  Provided by stmaryrd; we fall back to
%%% ASCII look-alikes if the package is not available.
\IfFileExists{stmaryrd.sty}{\usepackage{stmaryrd}}{%
  \providecommand{\llbracket}{[\![}%
  \providecommand{\rrbracket}{]\!]}%
}
\newcommand{\semof}[1]{\llbracket #1 \rrbracket}
\newcommand{\semofof}[2]{\semof{#1}(#2)}

%%% Up-to-scale entailment relation used for measurement-up-to-norm.
\newcommand{\entailedbyuptosc}{\mathrel{\models^{\mathit{uts}}}}


%%% =====================================================================
%%% 5. Transition-arrow rendering helpers
%%%
%%% Used in figures and inline transition diagrams.  Kept as-is from the
%%% legacy macros; they construct compact labelled arrows of the form
%%%       --[label]-->
%%% =====================================================================

\makeatletter
\DeclareRobustCommand{\shortto}{%
  \mathrel{\mathpalette\short@to\relax}%
}
\DeclareRobustCommand{\shortminus}{%
  \mathrel{\mathpalette\short@minus\relax}%
}
\newcommand{\short@to}[2]{%
  \mkern2mu
  \clipbox{{.5\width} 0 0 0}{$\m@th#1\vphantom{+}{\rightarrow}$}%
}
\newcommand{\short@minus}[2]{%
  \mkern2mu
  \clipbox{{.5\width} 0 0 0}{$\m@th#1\vphantom{+}{-}$}%
}
\makeatother

\newcommand{\labeledto}[1]{%
  \raisebox{-0.2pt}{\scalebox{1.2}{\ensuremath{%
    {\shortminus}\hspace{-1.7pt}%
    \raisebox{0.16ex}{$\scriptstyle[#1\hspace{-0.28pt}]$}%
    \hspace{-2.25pt}{\shortto}}}}}
\newcommand{\newlabeledto}[1]{%
  \mbox{$\longrightarrow\rlap{\hspace*{-4.95mm}%
    \raisebox{0.00mm}{\scalebox{0.8}{#1}}}$\!}}

\newcommand{\move}[1]{\labeledto{#1}}

%%% LTS-style transitions (used in transition systems, not LSTA proper).
\newcommand{\ltstr}{\mathit{Tr}}
\newcommand{\ltr}[1]{\xrightarrow{\smash{\raisebox{-0.7mm}{$\scriptstyle #1$}}}}

%%% Named-form transition macros from the legacy file (kept for the few
%%% places that build figures with them).
\newcommand{\trntransof}[4]{#1 \to #2(#3, #4)}
\newcommand{\transof}[5]{#1 \newlabeledto{#3} #2(#4, #5)}
\newcommand{\trnimageof}[2]{#1(#2)}
\newcommand{\trnimagestateof}[3]{\trnimageof{#1}{#2, #3}}


%%% =====================================================================
%%% 6. Quantum gates and operations
%%% =====================================================================

\newcommand{\ketof}[1]{\ket{#1}}

%%% Gate "wrappers" used in math mode.
\newcommand{\gateof}[1]{\mathrm{#1}}
\newcommand{\matrixof}[1]{\mathit{mat}(#1)}
\newcommand{\rowof}[1]{\mathit{row}(#1)}
\newcommand{\vectorof}[1]{\mathit{vec}(#1)}

%%% Common single- and two-qubit gates.
\newcommand{\had}{\mathsf{H}}
\newcommand{\hadtensor}{\had^{\otimes}}
\newcommand{\sgate}{\mathsf{S}}
\newcommand{\xgate}{\mathsf{X}}
\newcommand{\idgate}{\mathrm{Id}}
\newcommand{\cnot}{\mathsf{CX}}
\newcommand{\ccnot}{\mathsf{CCX}}
\newcommand{\rotx}{\mathsf{R}_{x}(\tfrac{\pi}{2})}

%%% Measurement.   \measurement{i}{=b}  -- measure qubit i, observe b.
\newcommand{\measurement}[2]{M_{#1}{#2}}

%%% Tree-transformation helpers (single-gate construction labels L/R).
\newcommand{\leftT}{\mathtt{L}}
\newcommand{\rightT}{\mathtt{R}}

%%% Classical store and qubit "directions" used in sLSTA constructions.
%%% These name the two principal directions of a classical branch.
\newcommand{\leftof}[1]{\mathtt{left}(#1)}
\newcommand{\rightof}[1]{\mathtt{right}(#1)}


%%% =====================================================================
%%% 7. Specifications and Hoare-style triples
%%% =====================================================================

%%% Standard Hoare triple {P} c {Q}.
\newcommand{\hoaretrip}[3]{\{#1\}~{#2}~\{#3\}}
%%% "Functional" triple with brackets (|P|) c (|Q|), used for specs
%%% intended to encode a Choi--Jamiolkowski-style state pair.
\newcommand{\assert}[1]{(\!|{#1}|\!)}
\newcommand{\funtrip}[3]{\assert{#1}\; {#2}\; \assert{#3}}

%%% Spec pair  c | psi >  used to bind a star pre/post to a state.
\newcommand{\specpair}[2]{#2\ket{#1}}
\newcommand{\specdef}[1]{#1\specpair{*}{}}

%%% Pre-/post-condition standard names (kept for cross-paper consistency).
\newcommand{\PreCond}{P}
\newcommand{\PostCond}{Q}
\newcommand{\pre}{\mathrm{Pre}}
\newcommand{\post}{\mathrm{Post}}

%%% Vectors of all-zeros / all-ones; the imaginary-zero subspace name.
\newcommand{\veczero}{\vec{0}}
\newcommand{\vecone}{\vec{1}}
\newcommand{\imzero}{\psi_{\Im}}


%%% =====================================================================
%%% 8. Tool names and method acronyms
%%%
%%% \tool refers to *our* tool; aliases below are for related-work and
%%% comparison tables.
%%% =====================================================================

\newcommand{\tool}{\textsc{CQ-AutoQ}\xspace}    % working name -- adjust later
\newcommand{\autoq}{\textsc{AutoQ}\xspace}
\newcommand{\sliqsim}{\textsc{SliQSim}\xspace}
\newcommand{\svsim}{\textsc{SV-Sim}\xspace}
\newcommand{\sliqec}{\textsc{SliQEC}\xspace}
\newcommand{\feynman}{\textsc{Feynman}\xspace}
\newcommand{\feynopt}{\textsc{Feynopt}\xspace}
\newcommand{\qcec}{\textsc{Qcec}\xspace}
\newcommand{\quasimodo}{\textsc{Quasimodo}\xspace}
\newcommand{\qbricks}{\textsc{Qbricks}\xspace}
\newcommand{\qiskit}{\textsc{Qiskit}\xspace}
\newcommand{\vata}{\textsc{Vata}\xspace}
\newcommand{\quartz}{\textsc{Quartz}\xspace}
\newcommand{\symqv}{\textsc{symQV}\xspace}
\newcommand{\caal}{\textsc{CaAL}\xspace}
\newcommand{\isabelle}{\textsc{Isabelle}\xspace}
\newcommand{\coq}{\textsc{Coq}\xspace}
\newcommand{\qasm}{\textsc{OpenQASM}\xspace}

%%% Method acronyms.
\newcommand{\lsta}{LSTA\xspace}
\newcommand{\lstas}{LSTAs\xspace}
\newcommand{\slsta}{sLSTA\xspace}
\newcommand{\slstas}{sLSTAs\xspace}

%%% Complexity classes (lower-case to avoid clash with cleveref labels).
\newcommand{\pspace}{\textsc{PSpace}\xspace}
\newcommand{\expspace}{\textsc{ExpSpace}\xspace}
\newcommand{\bigO}{\mathcal{O}}


%%% =====================================================================
%%% 9. Small tree-drawing helpers
%%%
%%% Inline TikZ snippets for compact tree literals in prose.  Pulled
%%% from the legacy file unchanged.
%%% =====================================================================

\newcommand{\smalltreeof}[3]{%
\scalebox{0.75}{%
\raisebox{3mm}{%
\tikz[level distance=2mm,sibling distance=8mm,baseline,inner sep=0.2mm]{%
\node {$#1$} child {node {$#2$}}%
child {node {$#3$}};}}}%
}

\newcommand{\bigtreeof}[6]{%
\scalebox{0.75}{%
\raisebox{3mm}{%
\tikz[level distance=2mm,
  level 1/.style={sibling distance=12mm},
  level 2/.style={sibling distance=8mm},
  baseline,inner sep=0.2mm]{%
\node {$#1$} child {node {$#2$}
                    child {node {$#3$}}
                    child {node {$#4$}}
                   }%
             child {node {$#2$}
                    child {node {$#5$}}
                    child {node {$#6$}}
                   };
}}}%
}


%%% =====================================================================
%%% 10. Miscellaneous
%%% =====================================================================

%%% Support set of a linear / weighted combination.
\newcommand{\supp}{\text{st}}
\newcommand{\suppof}[1]{\supp(#1)}

%%% Linear-form notation used in some proofs.
\newcommand{\linforms}{\mathbb{L}}
\newcommand{\linformsof}[1]{\linforms_{#1}}
\newcommand{\linof}[1]{\boldsymbol{#1}}
\newcommand{\linindexof}[2]{\linof{#1}[#2]}

%%% Positive part [x]_+.
\newcommand{\mm}[1]{[#1]_{+}}

%%% Misc symbols carried over from prior projects.
\newcommand{\sep}{\texttt{\#}}                % separator symbol
\newcommand{\initmark}{\to}
\newcommand{\partialto}{\rightharpoonup}
\newcommand{\compose}{\mathbin{\circ}}
\newcommand{\wildcard}{\mathtt{?}}

%%% Counting / encoding helpers.
\newcommand{\numof}[1]{\mathit{num}(#1)}
\newcommand{\revof}[1]{\mathit{rev}(#1)}
\newcommand{\encof}[1]{\mathit{enc}(#1)}

%%% A generic arithmetic operator placeholder used in some proofs.
\newcommand{\arithop}{\mathbin{\circledast}}
\newcommand{\minus}{\mathit{minus}}

%%% Choice / colour strings (legacy).
\newcommand{\colorstr}[1]{\mathbf{colorstr}(#1)}
\newcommand{\coloraut}[1]{\mathbf{dom}(#1)}

%%% ---- end of macros.tex ---------------------------------------------
   in main.tex.
%%%
%%% Conventions
%%%   - Sets in \mathbb           (\nat, \natz, \bools, \complex)
%%%   - Bit alphabet              \bit = \{0,1\}
%%%   - Tuples                    \tuple{a, b}
%%%   - Tool names                \autoq, \autq, \tool, ...
%%%   - LSTA / sLSTA names        \aut = \mathcal{A}; root states \rootstates
%%%   - Stem alphabet alias       \Store = Sigma_1 (rank-1 stem labels)
%%%   - Author notes              \ol, \yfc, \jal, \wl, \fil
%%%
%%% Sections
%%%   1.  Author notes and dev macros
%%%   2.  Sets, numbers, words
%%%   3.  Tree primitives (perfect binary, stemmed perfect binary)
%%%   4.  LSTA / sLSTA: states, transitions, runs, languages
%%%   5.  Transition-arrow rendering helpers
%%%   6.  Quantum gates and operations
%%%   7.  Specifications and Hoare-style triples
%%%   8.  Tool names and method acronyms
%%%   9.  Small tree-drawing helpers
%%%  10.  Miscellaneous
%%% =====================================================================


%%% =====================================================================
%%% 1. Author notes and dev macros
%%% =====================================================================

%%% Inline coloured author notes.  Toggle off by redefining to \gobble.
\newcommand{\ol}[1]{\textcolor{blue}{\ifmmode \text{[OL: #1]}\else [OL: #1] \fi}}
\newcommand{\yfc}[1]{\textcolor{green!50!blue}{\ifmmode \text{[YF: #1]}\else [YF: #1] \fi}}
\newcommand{\jal}[1]{\textcolor{pink}{\ifmmode \text{[JA: #1]}\else [JA: #1] \fi}}
\newcommand{\wl}[1]{\textcolor{orange}{\ifmmode \text{[WL: #1]}\else [WL: #1] \fi}}
\newcommand{\fil}[1]{\textcolor{red}{\ifmmode \text{[FI: #1]}\else [FI: #1] \fi}}

%%% Hide content (handy for parking paragraphs).  \showit kept for symmetry.
\newcommand{\hide}[1]{}
\newcommand{\showit}[1]{#1}

%%% Marker for places where a definition was added during a rewrite.
%%% Currently no-op; switch to the commented version to highlight.
% \newcommand{\newdef}[1]{\textcolor{purple}{[newdef: #1]}}
\newcommand{\newdef}[1]{}


%%% =====================================================================
%%% 2. Sets, numbers, words
%%% =====================================================================

\newcommand{\nat}{\mathbb{N}}                 % positive naturals (no 0)
\newcommand{\natz}{\nat_{0}}                  % naturals with 0
\newcommand{\naturals}{\nat}                  % alias
\newcommand{\integers}{\mathbb{Z}}
\newcommand{\bools}{\mathbb{B}}
\newcommand{\complex}{\mathbb{C}}
\newcommand{\cnum}{\complex}                  % older alias kept for back-compat

%%% Alphabets and bit strings
\newcommand{\bit}{\{0,1\}}                    % binary alphabet
\newcommand{\unary}{\{\mathtt{1}\}}           % unary alphabet (legacy)
\newcommand{\abc}{\Gamma}                     % generic alphabet
\newcommand{\colors}{\Omega}                  % colour / choice universe (legacy)

%%% The empty word.  Sometimes typeset epsilon, sometimes varepsilon; we
%%% standardise on \varepsilon for readability.
\newcommand{\emptyword}{\varepsilon}

%%% Tuples and scalable angle brackets.
%%% \tuple is a plain macro (matches existing usage q = \tuple{Q, ...}).
\newcommand{\tuple}[1]{\langle #1 \rangle}

%%% Common short fractions used in amplitudes.
\newcommand{\smallfrac}[2]{{\textstyle \frac{#1}{#2}}}
\newcommand{\invtwo}{\smallfrac{1}{2}}
\newcommand{\invsqrttwo}{\smallfrac{1}{\sqrt{2}}}


%%% =====================================================================
%%% 3. Tree primitives
%%%
%%% Covers both the perfect binary trees of [LSTA2025] and the
%%% stemmed perfect binary trees introduced in this paper.
%%% =====================================================================

%%% ---- generic tree language and domain --------------------------------
\newcommand{\trees}{\mathbb{T}}               % set of all (ranked) trees
\newcommand{\treesof}[1]{\trees_{#1}}
\newcommand{\abctrees}{\treesof{\abc}}

\newcommand{\dom}{\mathrm{dom}}               % \dom(T) -- domain of a tree
\newcommand{\domof}[1]{\dom(#1)}
\newcommand{\height}[1]{h(#1)}                % alternate height notation
\newcommand{\htof}[1]{\mathrm{ht}(#1)}        % preferred: \htof{v}

%%% Subtree, tree-cons, dummy zero tree, concatenation.
\newcommand{\subtreeof}[2]{#1_{|#2}}
\newcommand{\treecons}[3]{\mathsf{cons}(#1, #2, #3)}
\newcommand{\tzeroof}[1]{t^{#1}_{\mathit{zero}}}
\newcommand{\concat}{\mathbin{\bullet}}

%%% ---- ranked alphabets ------------------------------------------------
%%% Arity function on a ranked alphabet.  \arity{\Sigma} -> '#'.
%%% Defined as a macro so we can refer to it both in display ($\arity\colon
%%% \Sigma \to \{0,1,2\}$) and in prose ("\arity-1 symbols").
\newcommand{\arity}{\#}
\newcommand{\arityof}[1]{\arity(#1)}

%%% Stem-alphabet alias used by the sLSTA: Store = Sigma_1, the rank-1
%%% symbols carrying classical-store information.
\newcommand{\Store}{\mathtt{Store}}

%%% ---- shape of a stemmed perfect binary tree --------------------------
%%% (m, n) = (stem length, binary height).  Use \shape{m}{n}.
\newcommand{\shape}[2]{(#1, #2)}

%%% Root of a (possibly stemmed) tree T.  v_star with a star subscript.
\newcommand{\rootnode}{v_{\star}}


%%% =====================================================================
%%% 4. LSTA / sLSTA: states, transitions, runs, languages
%%% =====================================================================

%%% ---- automata names (calligraphic) ----------------------------------
\newcommand{\aut}{\mathcal{A}}
\newcommand{\autb}{\mathcal{B}}
\newcommand{\autm}{\mathcal{M}}
\newcommand{\autinv}{\mathcal{I}}
\newcommand{\cut}{\mathcal{C}}                % combined / context automaton
\newcommand{\nut}{\mathcal{N}}
\newcommand{\autex}{\aut_{\mathit{ex}}}       % a running-example automaton
\newcommand{\butex}{\autb_{\mathit{ex}}}
\newcommand{\autpre}{\aut_{\mathrm{pre}}}
\newcommand{\autpost}{\aut_{\mathrm{post}}}
\newcommand{\autbases}{\aut_{\mathit{bases}}}
\newcommand{\autRes}{\aut_{\mathsf{Result}}}
\newcommand{\auto}{\mathscr{A}}               % script-A (rarely used)

%%% ---- root states, transitions ----------------------------------------
\newcommand{\rootstates}{\mathcal{R}}

%%% A transition rendered as   q  --[C]-->  body.
%%% Body is either a binary symbol f(q0,q1), a unary symbol b(p), or a
%%% leaf symbol f.
\newcommand{\transtree}[3]{{#1}\xrightarrow{#2}~{#3}}

%%% Selectors on a transition delta.
\newcommand{\topof}[1]{\mathtt{top}(#1)}      % source state
\newcommand{\symof}[1]{\mathtt{sym}(#1)}      % label symbol
\newcommand{\botof}[1]{\mathtt{bot}(#1)}      % target state(s)

%%% ---- runs and level synchronization ---------------------------------
\newcommand{\run}{\rho}                       % canonical run variable
\newcommand{\level}{\mathrm{level}}           % \level(\run, d)
\newcommand{\choice}{\mathrm{choice}}         % \choice(\run)

%%% Language and choice-indexed language.
\newcommand{\lang}{\mathcal{L}}
\newcommand{\langof}[1]{\lang(#1)}

%%% Semantic brackets [[ . ]].  Provided by stmaryrd; we fall back to
%%% ASCII look-alikes if the package is not available.
\IfFileExists{stmaryrd.sty}{\usepackage{stmaryrd}}{%
  \providecommand{\llbracket}{[\![}%
  \providecommand{\rrbracket}{]\!]}%
}
\newcommand{\semof}[1]{\llbracket #1 \rrbracket}
\newcommand{\semofof}[2]{\semof{#1}(#2)}

%%% Up-to-scale entailment relation used for measurement-up-to-norm.
\newcommand{\entailedbyuptosc}{\mathrel{\models^{\mathit{uts}}}}


%%% =====================================================================
%%% 5. Transition-arrow rendering helpers
%%%
%%% Used in figures and inline transition diagrams.  Kept as-is from the
%%% legacy macros; they construct compact labelled arrows of the form
%%%       --[label]-->
%%% =====================================================================

\makeatletter
\DeclareRobustCommand{\shortto}{%
  \mathrel{\mathpalette\short@to\relax}%
}
\DeclareRobustCommand{\shortminus}{%
  \mathrel{\mathpalette\short@minus\relax}%
}
\newcommand{\short@to}[2]{%
  \mkern2mu
  \clipbox{{.5\width} 0 0 0}{$\m@th#1\vphantom{+}{\rightarrow}$}%
}
\newcommand{\short@minus}[2]{%
  \mkern2mu
  \clipbox{{.5\width} 0 0 0}{$\m@th#1\vphantom{+}{-}$}%
}
\makeatother

\newcommand{\labeledto}[1]{%
  \raisebox{-0.2pt}{\scalebox{1.2}{\ensuremath{%
    {\shortminus}\hspace{-1.7pt}%
    \raisebox{0.16ex}{$\scriptstyle[#1\hspace{-0.28pt}]$}%
    \hspace{-2.25pt}{\shortto}}}}}
\newcommand{\newlabeledto}[1]{%
  \mbox{$\longrightarrow\rlap{\hspace*{-4.95mm}%
    \raisebox{0.00mm}{\scalebox{0.8}{#1}}}$\!}}

\newcommand{\move}[1]{\labeledto{#1}}

%%% LTS-style transitions (used in transition systems, not LSTA proper).
\newcommand{\ltstr}{\mathit{Tr}}
\newcommand{\ltr}[1]{\xrightarrow{\smash{\raisebox{-0.7mm}{$\scriptstyle #1$}}}}

%%% Named-form transition macros from the legacy file (kept for the few
%%% places that build figures with them).
\newcommand{\trntransof}[4]{#1 \to #2(#3, #4)}
\newcommand{\transof}[5]{#1 \newlabeledto{#3} #2(#4, #5)}
\newcommand{\trnimageof}[2]{#1(#2)}
\newcommand{\trnimagestateof}[3]{\trnimageof{#1}{#2, #3}}


%%% =====================================================================
%%% 6. Quantum gates and operations
%%% =====================================================================

\newcommand{\ketof}[1]{\ket{#1}}

%%% Gate "wrappers" used in math mode.
\newcommand{\gateof}[1]{\mathrm{#1}}
\newcommand{\matrixof}[1]{\mathit{mat}(#1)}
\newcommand{\rowof}[1]{\mathit{row}(#1)}
\newcommand{\vectorof}[1]{\mathit{vec}(#1)}

%%% Common single- and two-qubit gates.
\newcommand{\had}{\mathsf{H}}
\newcommand{\hadtensor}{\had^{\otimes}}
\newcommand{\sgate}{\mathsf{S}}
\newcommand{\xgate}{\mathsf{X}}
\newcommand{\idgate}{\mathrm{Id}}
\newcommand{\cnot}{\mathsf{CX}}
\newcommand{\ccnot}{\mathsf{CCX}}
\newcommand{\rotx}{\mathsf{R}_{x}(\tfrac{\pi}{2})}

%%% Measurement.   \measurement{i}{=b}  -- measure qubit i, observe b.
\newcommand{\measurement}[2]{M_{#1}{#2}}

%%% Tree-transformation helpers (single-gate construction labels L/R).
\newcommand{\leftT}{\mathtt{L}}
\newcommand{\rightT}{\mathtt{R}}

%%% Classical store and qubit "directions" used in sLSTA constructions.
%%% These name the two principal directions of a classical branch.
\newcommand{\leftof}[1]{\mathtt{left}(#1)}
\newcommand{\rightof}[1]{\mathtt{right}(#1)}


%%% =====================================================================
%%% 7. Specifications and Hoare-style triples
%%% =====================================================================

%%% Standard Hoare triple {P} c {Q}.
\newcommand{\hoaretrip}[3]{\{#1\}~{#2}~\{#3\}}
%%% "Functional" triple with brackets (|P|) c (|Q|), used for specs
%%% intended to encode a Choi--Jamiolkowski-style state pair.
\newcommand{\assert}[1]{(\!|{#1}|\!)}
\newcommand{\funtrip}[3]{\assert{#1}\; {#2}\; \assert{#3}}

%%% Spec pair  c | psi >  used to bind a star pre/post to a state.
\newcommand{\specpair}[2]{#2\ket{#1}}
\newcommand{\specdef}[1]{#1\specpair{*}{}}

%%% Pre-/post-condition standard names (kept for cross-paper consistency).
\newcommand{\PreCond}{P}
\newcommand{\PostCond}{Q}
\newcommand{\pre}{\mathrm{Pre}}
\newcommand{\post}{\mathrm{Post}}

%%% Vectors of all-zeros / all-ones; the imaginary-zero subspace name.
\newcommand{\veczero}{\vec{0}}
\newcommand{\vecone}{\vec{1}}
\newcommand{\imzero}{\psi_{\Im}}


%%% =====================================================================
%%% 8. Tool names and method acronyms
%%%
%%% \tool refers to *our* tool; aliases below are for related-work and
%%% comparison tables.
%%% =====================================================================

\newcommand{\tool}{\textsc{CQ-AutoQ}\xspace}    % working name -- adjust later
\newcommand{\autoq}{\textsc{AutoQ}\xspace}
\newcommand{\sliqsim}{\textsc{SliQSim}\xspace}
\newcommand{\svsim}{\textsc{SV-Sim}\xspace}
\newcommand{\sliqec}{\textsc{SliQEC}\xspace}
\newcommand{\feynman}{\textsc{Feynman}\xspace}
\newcommand{\feynopt}{\textsc{Feynopt}\xspace}
\newcommand{\qcec}{\textsc{Qcec}\xspace}
\newcommand{\quasimodo}{\textsc{Quasimodo}\xspace}
\newcommand{\qbricks}{\textsc{Qbricks}\xspace}
\newcommand{\qiskit}{\textsc{Qiskit}\xspace}
\newcommand{\vata}{\textsc{Vata}\xspace}
\newcommand{\quartz}{\textsc{Quartz}\xspace}
\newcommand{\symqv}{\textsc{symQV}\xspace}
\newcommand{\caal}{\textsc{CaAL}\xspace}
\newcommand{\isabelle}{\textsc{Isabelle}\xspace}
\newcommand{\coq}{\textsc{Coq}\xspace}
\newcommand{\qasm}{\textsc{OpenQASM}\xspace}

%%% Method acronyms.
\newcommand{\lsta}{LSTA\xspace}
\newcommand{\lstas}{LSTAs\xspace}
\newcommand{\slsta}{sLSTA\xspace}
\newcommand{\slstas}{sLSTAs\xspace}

%%% Complexity classes (lower-case to avoid clash with cleveref labels).
\newcommand{\pspace}{\textsc{PSpace}\xspace}
\newcommand{\expspace}{\textsc{ExpSpace}\xspace}
\newcommand{\bigO}{\mathcal{O}}


%%% =====================================================================
%%% 9. Small tree-drawing helpers
%%%
%%% Inline TikZ snippets for compact tree literals in prose.  Pulled
%%% from the legacy file unchanged.
%%% =====================================================================

\newcommand{\smalltreeof}[3]{%
\scalebox{0.75}{%
\raisebox{3mm}{%
\tikz[level distance=2mm,sibling distance=8mm,baseline,inner sep=0.2mm]{%
\node {$#1$} child {node {$#2$}}%
child {node {$#3$}};}}}%
}

\newcommand{\bigtreeof}[6]{%
\scalebox{0.75}{%
\raisebox{3mm}{%
\tikz[level distance=2mm,
  level 1/.style={sibling distance=12mm},
  level 2/.style={sibling distance=8mm},
  baseline,inner sep=0.2mm]{%
\node {$#1$} child {node {$#2$}
                    child {node {$#3$}}
                    child {node {$#4$}}
                   }%
             child {node {$#2$}
                    child {node {$#5$}}
                    child {node {$#6$}}
                   };
}}}%
}


%%% =====================================================================
%%% 10. Miscellaneous
%%% =====================================================================

%%% Support set of a linear / weighted combination.
\newcommand{\supp}{\text{st}}
\newcommand{\suppof}[1]{\supp(#1)}

%%% Linear-form notation used in some proofs.
\newcommand{\linforms}{\mathbb{L}}
\newcommand{\linformsof}[1]{\linforms_{#1}}
\newcommand{\linof}[1]{\boldsymbol{#1}}
\newcommand{\linindexof}[2]{\linof{#1}[#2]}

%%% Positive part [x]_+.
\newcommand{\mm}[1]{[#1]_{+}}

%%% Misc symbols carried over from prior projects.
\newcommand{\sep}{\texttt{\#}}                % separator symbol
\newcommand{\initmark}{\to}
\newcommand{\partialto}{\rightharpoonup}
\newcommand{\compose}{\mathbin{\circ}}
\newcommand{\wildcard}{\mathtt{?}}

%%% Counting / encoding helpers.
\newcommand{\numof}[1]{\mathit{num}(#1)}
\newcommand{\revof}[1]{\mathit{rev}(#1)}
\newcommand{\encof}[1]{\mathit{enc}(#1)}

%%% A generic arithmetic operator placeholder used in some proofs.
\newcommand{\arithop}{\mathbin{\circledast}}
\newcommand{\minus}{\mathit{minus}}

%%% Choice / colour strings (legacy).
\newcommand{\colorstr}[1]{\mathbf{colorstr}(#1)}
\newcommand{\coloraut}[1]{\mathbf{dom}(#1)}

%%% ---- end of macros.tex ---------------------------------------------
   in main.tex.
%%%
%%% Conventions
%%%   - Sets in \mathbb           (\nat, \natz, \bools, \complex)
%%%   - Bit alphabet              \bit = \{0,1\}
%%%   - Tuples                    \tuple{a, b}
%%%   - Tool names                \autoq, \autq, \tool, ...
%%%   - LSTA / sLSTA names        \aut = \mathcal{A}; root states \rootstates
%%%   - Stem alphabet alias       \Store = Sigma_1 (rank-1 stem labels)
%%%   - Author notes              \ol, \yfc, \jal, \wl, \fil
%%%
%%% Sections
%%%   1.  Author notes and dev macros
%%%   2.  Sets, numbers, words
%%%   3.  Tree primitives (perfect binary, stemmed perfect binary)
%%%   4.  LSTA / sLSTA: states, transitions, runs, languages
%%%   5.  Transition-arrow rendering helpers
%%%   6.  Quantum gates and operations
%%%   7.  Specifications and Hoare-style triples
%%%   8.  Tool names and method acronyms
%%%   9.  Small tree-drawing helpers
%%%  10.  Miscellaneous
%%% =====================================================================


%%% =====================================================================
%%% 1. Author notes and dev macros
%%% =====================================================================

%%% Inline coloured author notes.  Toggle off by redefining to \gobble.
\newcommand{\ol}[1]{\textcolor{blue}{\ifmmode \text{[OL: #1]}\else [OL: #1] \fi}}
\newcommand{\yfc}[1]{\textcolor{green!50!blue}{\ifmmode \text{[YF: #1]}\else [YF: #1] \fi}}
\newcommand{\jal}[1]{\textcolor{pink}{\ifmmode \text{[JA: #1]}\else [JA: #1] \fi}}
\newcommand{\wl}[1]{\textcolor{orange}{\ifmmode \text{[WL: #1]}\else [WL: #1] \fi}}
\newcommand{\fil}[1]{\textcolor{red}{\ifmmode \text{[FI: #1]}\else [FI: #1] \fi}}

%%% Hide content (handy for parking paragraphs).  \showit kept for symmetry.
\newcommand{\hide}[1]{}
\newcommand{\showit}[1]{#1}

%%% Marker for places where a definition was added during a rewrite.
%%% Currently no-op; switch to the commented version to highlight.
% \newcommand{\newdef}[1]{\textcolor{purple}{[newdef: #1]}}
\newcommand{\newdef}[1]{}


%%% =====================================================================
%%% 2. Sets, numbers, words
%%% =====================================================================

\newcommand{\nat}{\mathbb{N}}                 % positive naturals (no 0)
\newcommand{\natz}{\nat_{0}}                  % naturals with 0
\newcommand{\naturals}{\nat}                  % alias
\newcommand{\integers}{\mathbb{Z}}
\newcommand{\bools}{\mathbb{B}}
\newcommand{\complex}{\mathbb{C}}
\newcommand{\cnum}{\complex}                  % older alias kept for back-compat

%%% Alphabets and bit strings
\newcommand{\bit}{\{0,1\}}                    % binary alphabet
\newcommand{\unary}{\{\mathtt{1}\}}           % unary alphabet (legacy)
\newcommand{\abc}{\Gamma}                     % generic alphabet
\newcommand{\colors}{\Omega}                  % colour / choice universe (legacy)

%%% The empty word.  Sometimes typeset epsilon, sometimes varepsilon; we
%%% standardise on \varepsilon for readability.
\newcommand{\emptyword}{\varepsilon}

%%% Tuples and scalable angle brackets.
%%% \tuple is a plain macro (matches existing usage q = \tuple{Q, ...}).
\newcommand{\tuple}[1]{\langle #1 \rangle}

%%% Common short fractions used in amplitudes.
\newcommand{\smallfrac}[2]{{\textstyle \frac{#1}{#2}}}
\newcommand{\invtwo}{\smallfrac{1}{2}}
\newcommand{\invsqrttwo}{\smallfrac{1}{\sqrt{2}}}


%%% =====================================================================
%%% 3. Tree primitives
%%%
%%% Covers both the perfect binary trees of [LSTA2025] and the
%%% stemmed perfect binary trees introduced in this paper.
%%% =====================================================================

%%% ---- generic tree language and domain --------------------------------
\newcommand{\trees}{\mathbb{T}}               % set of all (ranked) trees
\newcommand{\treesof}[1]{\trees_{#1}}
\newcommand{\abctrees}{\treesof{\abc}}

\newcommand{\dom}{\mathrm{dom}}               % \dom(T) -- domain of a tree
\newcommand{\domof}[1]{\dom(#1)}
\newcommand{\height}[1]{h(#1)}                % alternate height notation
\newcommand{\htof}[1]{\mathrm{ht}(#1)}        % preferred: \htof{v}

%%% Subtree, tree-cons, dummy zero tree, concatenation.
\newcommand{\subtreeof}[2]{#1_{|#2}}
\newcommand{\treecons}[3]{\mathsf{cons}(#1, #2, #3)}
\newcommand{\tzeroof}[1]{t^{#1}_{\mathit{zero}}}
\newcommand{\concat}{\mathbin{\bullet}}

%%% ---- ranked alphabets ------------------------------------------------
%%% Arity function on a ranked alphabet.  \arity{\Sigma} -> '#'.
%%% Defined as a macro so we can refer to it both in display ($\arity\colon
%%% \Sigma \to \{0,1,2\}$) and in prose ("\arity-1 symbols").
\newcommand{\arity}{\#}
\newcommand{\arityof}[1]{\arity(#1)}

%%% Stem-alphabet alias used by the sLSTA: Store = Sigma_1, the rank-1
%%% symbols carrying classical-store information.
\newcommand{\Store}{\mathtt{Store}}

%%% ---- shape of a stemmed perfect binary tree --------------------------
%%% (m, n) = (stem length, binary height).  Use \shape{m}{n}.
\newcommand{\shape}[2]{(#1, #2)}

%%% Root of a (possibly stemmed) tree T.  v_star with a star subscript.
\newcommand{\rootnode}{v_{\star}}


%%% =====================================================================
%%% 4. LSTA / sLSTA: states, transitions, runs, languages
%%% =====================================================================

%%% ---- automata names (calligraphic) ----------------------------------
\newcommand{\aut}{\mathcal{A}}
\newcommand{\autb}{\mathcal{B}}
\newcommand{\autm}{\mathcal{M}}
\newcommand{\autinv}{\mathcal{I}}
\newcommand{\cut}{\mathcal{C}}                % combined / context automaton
\newcommand{\nut}{\mathcal{N}}
\newcommand{\autex}{\aut_{\mathit{ex}}}       % a running-example automaton
\newcommand{\butex}{\autb_{\mathit{ex}}}
\newcommand{\autpre}{\aut_{\mathrm{pre}}}
\newcommand{\autpost}{\aut_{\mathrm{post}}}
\newcommand{\autbases}{\aut_{\mathit{bases}}}
\newcommand{\autRes}{\aut_{\mathsf{Result}}}
\newcommand{\auto}{\mathscr{A}}               % script-A (rarely used)

%%% ---- root states, transitions ----------------------------------------
\newcommand{\rootstates}{\mathcal{R}}

%%% A transition rendered as   q  --[C]-->  body.
%%% Body is either a binary symbol f(q0,q1), a unary symbol b(p), or a
%%% leaf symbol f.
\newcommand{\transtree}[3]{{#1}\xrightarrow{#2}~{#3}}

%%% Selectors on a transition delta.
\newcommand{\topof}[1]{\mathtt{top}(#1)}      % source state
\newcommand{\symof}[1]{\mathtt{sym}(#1)}      % label symbol
\newcommand{\botof}[1]{\mathtt{bot}(#1)}      % target state(s)

%%% ---- runs and level synchronization ---------------------------------
\newcommand{\run}{\rho}                       % canonical run variable
\newcommand{\level}{\mathrm{level}}           % \level(\run, d)
\newcommand{\choice}{\mathrm{choice}}         % \choice(\run)

%%% Language and choice-indexed language.
\newcommand{\lang}{\mathcal{L}}
\newcommand{\langof}[1]{\lang(#1)}

%%% Semantic brackets [[ . ]].  Provided by stmaryrd; we fall back to
%%% ASCII look-alikes if the package is not available.
\IfFileExists{stmaryrd.sty}{\usepackage{stmaryrd}}{%
  \providecommand{\llbracket}{[\![}%
  \providecommand{\rrbracket}{]\!]}%
}
\newcommand{\semof}[1]{\llbracket #1 \rrbracket}
\newcommand{\semofof}[2]{\semof{#1}(#2)}

%%% Up-to-scale entailment relation used for measurement-up-to-norm.
\newcommand{\entailedbyuptosc}{\mathrel{\models^{\mathit{uts}}}}


%%% =====================================================================
%%% 5. Transition-arrow rendering helpers
%%%
%%% Used in figures and inline transition diagrams.  Kept as-is from the
%%% legacy macros; they construct compact labelled arrows of the form
%%%       --[label]-->
%%% =====================================================================

\makeatletter
\DeclareRobustCommand{\shortto}{%
  \mathrel{\mathpalette\short@to\relax}%
}
\DeclareRobustCommand{\shortminus}{%
  \mathrel{\mathpalette\short@minus\relax}%
}
\newcommand{\short@to}[2]{%
  \mkern2mu
  \clipbox{{.5\width} 0 0 0}{$\m@th#1\vphantom{+}{\rightarrow}$}%
}
\newcommand{\short@minus}[2]{%
  \mkern2mu
  \clipbox{{.5\width} 0 0 0}{$\m@th#1\vphantom{+}{-}$}%
}
\makeatother

\newcommand{\labeledto}[1]{%
  \raisebox{-0.2pt}{\scalebox{1.2}{\ensuremath{%
    {\shortminus}\hspace{-1.7pt}%
    \raisebox{0.16ex}{$\scriptstyle[#1\hspace{-0.28pt}]$}%
    \hspace{-2.25pt}{\shortto}}}}}
\newcommand{\newlabeledto}[1]{%
  \mbox{$\longrightarrow\rlap{\hspace*{-4.95mm}%
    \raisebox{0.00mm}{\scalebox{0.8}{#1}}}$\!}}

\newcommand{\move}[1]{\labeledto{#1}}

%%% LTS-style transitions (used in transition systems, not LSTA proper).
\newcommand{\ltstr}{\mathit{Tr}}
\newcommand{\ltr}[1]{\xrightarrow{\smash{\raisebox{-0.7mm}{$\scriptstyle #1$}}}}

%%% Named-form transition macros from the legacy file (kept for the few
%%% places that build figures with them).
\newcommand{\trntransof}[4]{#1 \to #2(#3, #4)}
\newcommand{\transof}[5]{#1 \newlabeledto{#3} #2(#4, #5)}
\newcommand{\trnimageof}[2]{#1(#2)}
\newcommand{\trnimagestateof}[3]{\trnimageof{#1}{#2, #3}}


%%% =====================================================================
%%% 6. Quantum gates and operations
%%% =====================================================================

\newcommand{\ketof}[1]{\ket{#1}}

%%% Gate "wrappers" used in math mode.
\newcommand{\gateof}[1]{\mathrm{#1}}
\newcommand{\matrixof}[1]{\mathit{mat}(#1)}
\newcommand{\rowof}[1]{\mathit{row}(#1)}
\newcommand{\vectorof}[1]{\mathit{vec}(#1)}

%%% Common single- and two-qubit gates.
\newcommand{\had}{\mathsf{H}}
\newcommand{\hadtensor}{\had^{\otimes}}
\newcommand{\sgate}{\mathsf{S}}
\newcommand{\xgate}{\mathsf{X}}
\newcommand{\idgate}{\mathrm{Id}}
\newcommand{\cnot}{\mathsf{CX}}
\newcommand{\ccnot}{\mathsf{CCX}}
\newcommand{\rotx}{\mathsf{R}_{x}(\tfrac{\pi}{2})}

%%% Measurement.   \measurement{i}{=b}  -- measure qubit i, observe b.
\newcommand{\measurement}[2]{M_{#1}{#2}}

%%% Tree-transformation helpers (single-gate construction labels L/R).
\newcommand{\leftT}{\mathtt{L}}
\newcommand{\rightT}{\mathtt{R}}

%%% Classical store and qubit "directions" used in sLSTA constructions.
%%% These name the two principal directions of a classical branch.
\newcommand{\leftof}[1]{\mathtt{left}(#1)}
\newcommand{\rightof}[1]{\mathtt{right}(#1)}


%%% =====================================================================
%%% 7. Specifications and Hoare-style triples
%%% =====================================================================

%%% Standard Hoare triple {P} c {Q}.
\newcommand{\hoaretrip}[3]{\{#1\}~{#2}~\{#3\}}
%%% "Functional" triple with brackets (|P|) c (|Q|), used for specs
%%% intended to encode a Choi--Jamiolkowski-style state pair.
\newcommand{\assert}[1]{(\!|{#1}|\!)}
\newcommand{\funtrip}[3]{\assert{#1}\; {#2}\; \assert{#3}}

%%% Spec pair  c | psi >  used to bind a star pre/post to a state.
\newcommand{\specpair}[2]{#2\ket{#1}}
\newcommand{\specdef}[1]{#1\specpair{*}{}}

%%% Pre-/post-condition standard names (kept for cross-paper consistency).
\newcommand{\PreCond}{P}
\newcommand{\PostCond}{Q}
\newcommand{\pre}{\mathrm{Pre}}
\newcommand{\post}{\mathrm{Post}}

%%% Vectors of all-zeros / all-ones; the imaginary-zero subspace name.
\newcommand{\veczero}{\vec{0}}
\newcommand{\vecone}{\vec{1}}
\newcommand{\imzero}{\psi_{\Im}}


%%% =====================================================================
%%% 8. Tool names and method acronyms
%%%
%%% \tool refers to *our* tool; aliases below are for related-work and
%%% comparison tables.
%%% =====================================================================

\newcommand{\tool}{\textsc{CQ-AutoQ}\xspace}    % working name -- adjust later
\newcommand{\autoq}{\textsc{AutoQ}\xspace}
\newcommand{\sliqsim}{\textsc{SliQSim}\xspace}
\newcommand{\svsim}{\textsc{SV-Sim}\xspace}
\newcommand{\sliqec}{\textsc{SliQEC}\xspace}
\newcommand{\feynman}{\textsc{Feynman}\xspace}
\newcommand{\feynopt}{\textsc{Feynopt}\xspace}
\newcommand{\qcec}{\textsc{Qcec}\xspace}
\newcommand{\quasimodo}{\textsc{Quasimodo}\xspace}
\newcommand{\qbricks}{\textsc{Qbricks}\xspace}
\newcommand{\qiskit}{\textsc{Qiskit}\xspace}
\newcommand{\vata}{\textsc{Vata}\xspace}
\newcommand{\quartz}{\textsc{Quartz}\xspace}
\newcommand{\symqv}{\textsc{symQV}\xspace}
\newcommand{\caal}{\textsc{CaAL}\xspace}
\newcommand{\isabelle}{\textsc{Isabelle}\xspace}
\newcommand{\coq}{\textsc{Coq}\xspace}
\newcommand{\qasm}{\textsc{OpenQASM}\xspace}

%%% Method acronyms.
\newcommand{\lsta}{LSTA\xspace}
\newcommand{\lstas}{LSTAs\xspace}
\newcommand{\slsta}{sLSTA\xspace}
\newcommand{\slstas}{sLSTAs\xspace}

%%% Complexity classes (lower-case to avoid clash with cleveref labels).
\newcommand{\pspace}{\textsc{PSpace}\xspace}
\newcommand{\expspace}{\textsc{ExpSpace}\xspace}
\newcommand{\bigO}{\mathcal{O}}


%%% =====================================================================
%%% 9. Small tree-drawing helpers
%%%
%%% Inline TikZ snippets for compact tree literals in prose.  Pulled
%%% from the legacy file unchanged.
%%% =====================================================================

\newcommand{\smalltreeof}[3]{%
\scalebox{0.75}{%
\raisebox{3mm}{%
\tikz[level distance=2mm,sibling distance=8mm,baseline,inner sep=0.2mm]{%
\node {$#1$} child {node {$#2$}}%
child {node {$#3$}};}}}%
}

\newcommand{\bigtreeof}[6]{%
\scalebox{0.75}{%
\raisebox{3mm}{%
\tikz[level distance=2mm,
  level 1/.style={sibling distance=12mm},
  level 2/.style={sibling distance=8mm},
  baseline,inner sep=0.2mm]{%
\node {$#1$} child {node {$#2$}
                    child {node {$#3$}}
                    child {node {$#4$}}
                   }%
             child {node {$#2$}
                    child {node {$#5$}}
                    child {node {$#6$}}
                   };
}}}%
}


%%% =====================================================================
%%% 10. Miscellaneous
%%% =====================================================================

%%% Support set of a linear / weighted combination.
\newcommand{\supp}{\text{st}}
\newcommand{\suppof}[1]{\supp(#1)}

%%% Linear-form notation used in some proofs.
\newcommand{\linforms}{\mathbb{L}}
\newcommand{\linformsof}[1]{\linforms_{#1}}
\newcommand{\linof}[1]{\boldsymbol{#1}}
\newcommand{\linindexof}[2]{\linof{#1}[#2]}

%%% Positive part [x]_+.
\newcommand{\mm}[1]{[#1]_{+}}

%%% Misc symbols carried over from prior projects.
\newcommand{\sep}{\texttt{\#}}                % separator symbol
\newcommand{\initmark}{\to}
\newcommand{\partialto}{\rightharpoonup}
\newcommand{\compose}{\mathbin{\circ}}
\newcommand{\wildcard}{\mathtt{?}}

%%% Counting / encoding helpers.
\newcommand{\numof}[1]{\mathit{num}(#1)}
\newcommand{\revof}[1]{\mathit{rev}(#1)}
\newcommand{\encof}[1]{\mathit{enc}(#1)}

%%% A generic arithmetic operator placeholder used in some proofs.
\newcommand{\arithop}{\mathbin{\circledast}}
\newcommand{\minus}{\mathit{minus}}

%%% Choice / colour strings (legacy).
\newcommand{\colorstr}[1]{\mathbf{colorstr}(#1)}
\newcommand{\coloraut}[1]{\mathbf{dom}(#1)}

%%% ---- end of macros.tex ---------------------------------------------
